% ==============================================================================
% CHAPTER 4: KEY TECHNOLOGIES IN MULTI-AGENT COLLABORATIVE QA
% ==============================================================================
\section{Key Technologies in Multi-Agent Collaborative QA}
\label{sec:key_tech}

The efficient operation of a multi-agent collaborative question answering (QA) system depends not only on a well-designed organizational architecture but also on the support of core underlying technologies.
This section investigates four technological dimensions: communication mechanisms, task decomposition, memory management, and reasoning enhancement.

% ── 4.1  Communication Mechanisms ─────────────────────────────────────────────
\subsection{Agent Communication Mechanisms}
\label{sec:comm}

\subsubsection{Explicit Communication}
\label{sec:explicit_comm}

Explicit communication refers to the direct transfer of information between agents via predefined formats.
The most common form is message passing based on structured languages (JSON,
XML) or dedicated prompt templates.
AutoGen \citep{wu2024autogen} provides a representative implementation:
customisable conversable agents send and receive natural-language or code
messages following conversation patterns defined by developers.
ChatDev \citep{qian2023chatdev}, by contrast, organises role-specialised agents
through a multi-turn dialogue chain covering design, coding, and testing stages.


\subsubsection{Implicit Communication}
\label{sec:implicit_comm}

Implicit communication does not require direct agent-to-agent messaging; instead, collaboration is mediated via modifications to a shared operating environment.
For instance, in a code-generation QA task, one agent alters the state of a code repository or database, while another agent obtains feedback by observing compilation errors or test execution results.

This paradigm can reduce redundant message transfers, though its impact on communication overhead and fault tolerance is contingent on the implementation and task environment.

\subsubsection{Standardized Communication Protocols}
\label{sec:protocols}
With the rapid development of multi-agent ecosystems, interoperability barriers have arisen between agents developed by different vendors.
To address this issue, the industry has initiated efforts to standardize communication protocols. For example, Anthropic introduced the Model Context Protocol (MCP)
\citep{anthropic2024mcp}—dubbed the ``USB-C of AI''—to standardise
the way agents connect to external tools and data sources.
Google proposed the Agent-to-Agent (A2A) protocol \citep{google2025a2a} to
regulate horizontal peer-to-peer collaboration among agents.
These standards are being gradually adopted by the open-source
community and the Linux Foundation, laying a foundation for heterogeneous
multi-agent QA systems.

% ── 4.2  Task Decomposition and Scheduling ────────────────────────────────────
\subsection{Task Decomposition and Scheduling}
\label{sec:task_decomp}

\subsubsection{Hierarchical Task Decomposition}
\label{sec:hier_decomp}

For complex QA tasks, LLMs often cannot answer in a single step.
Hierarchical task decomposition reduces difficulty by breaking complex queries
into sub-tasks that can be processed in parallel or sequentially.
Least-to-Most Prompting \citep{zhou2023least_to_most} and Decomposed
Prompting \citep{khot2023decomposed_prompting} allow a model to solve simpler
sub-problems first and then exploit their answers to resolve the ultimate
target; Plan-and-Solve \citep{wang2023plan_and_solve} explicitly separates
the planning and solving phases, substantially improving decomposition
reliability.
Formally, given a complex query $Q$, a decomposition function $\mathcal{D}$
maps it to an ordered sequence of sub-tasks:
\begin{equation}
  \begin{aligned}
    \mathcal{D}(Q) &\to \{q_1, q_2, \ldots, q_m\}, \\
    \operatorname{Answer}(Q) &=
    \mathcal{F}\!\left(\operatorname{Sol}(q_1),\ldots,\operatorname{Sol}(q_m)\right)
  \end{aligned}
  \label{eq:decompose}
\end{equation}
where $\operatorname{Sol}(q_i)$ is a sub-agent's solution to sub-task $q_i$
and $\mathcal{F}$ is an aggregation function (e.g., voting, concatenation, or
re-generation).

\subsubsection{Role-Based Task Assignment}
\label{sec:role_assign}

After decomposition, the system must assign sub-tasks to the most suitable
agents.
MetaGPT \citep{hong2024metagpt} and AgentVerse \citep{chen2024agentverse}
construct fine-grained agent profiles (Agent Profiling) that tag each agent
with specific role labels.
During runtime, a scheduler uses skill-based routing to dispatch sub-tasks
such as code writing, literature retrieval, or logical verification to the
corresponding expert agent, maximising throughput.

\subsubsection{Pipeline Optimization and Adaptive Orchestration}
\label{sec:load_balance}

Static assignment can be inefficient when task difficulty and agent capabilities vary across queries. DSPy \citep{khattab2024dspy} addresses a related but narrower problem: it represents language-model pipelines as declarative modules and compiles their prompts and demonstrations against a task metric. It should therefore be regarded as a pipeline-optimization toolkit rather than evidence supporting runtime multi-agent load balancing. Dynamic agent recruitment, routing, and early stopping remain distinct orchestration decisions that require task-specific evaluation.

% ── 4.3  Memory and Knowledge Management ──────────────────────────────────────
\subsection{Memory and Knowledge Management}
\label{sec:memory}

\subsubsection{Short-Term Memory and Working Memory}
\label{sec:short_mem}

Short-term memory preserves dialogue history and local context for the
current QA session.
To prevent the model from forgetting in extended conversations, MemGPT
\citep{packer2023memgpt} draws on the paged-memory mechanism of modern
operating systems to design a hierarchical working-memory system.
The system retains critical context in memory based on attention weights
and evicts redundant information, providing seamless long-conversation
support within a limited context window.

\subsubsection{Long-Term Memory and Vector Retrieval (RAG)}
\label{sec:long_mem}

Long-term memory typically relies on vector databases to store vast
historical experiences and external domain knowledge
\citep{zhang2024memory_survey}.
Retrieval-Augmented Generation (RAG) \citep{lewis2020rag,gao2024rag_survey}
is the primary mechanism through which multi-agent QA systems access
long-term memory.
A retrieval agent selects the Top-$k$ relevant documents from knowledge
base $\mathcal{D}$ using cosine similarity:
\begin{equation}
  \operatorname{Rel}(q,d_i) =
  \frac{\boldsymbol{\phi}(q)^{\top}\boldsymbol{\phi}(d_i)}
       {\|\boldsymbol{\phi}(q)\|\,\|\boldsymbol{\phi}(d_i)\|},\qquad
  \hat{\mathcal{D}} =
  \operatorname*{Top\text{-}k}_{d_i\in\mathcal{D}}\operatorname{Rel}(q,d_i)
  \label{eq:rag}
\end{equation}
where $\boldsymbol{\phi}(\cdot)$ is the dense encoder embedding function and
$\hat{\mathcal{D}}$ is the set of retrieved relevant documents.
Retrieval has evolved from early passive Naive RAG to Self-RAG
\citep{asai2024selfrag} and proactive retrieval (FLARE)
\citep{jiang2023active_rag}; specialised retrieval agents can now
autonomously determine when to retrieve, evaluate retrieval quality, and
pass refined knowledge to the generation agent.

\subsubsection{Cross-Agent Knowledge Sharing}
\label{sec:knowledge_share}

In complex collaboration, knowledge must not be siloed within individual
agents.
The system needs a global knowledge graph shared across agents so that all
collaborative nodes can synchronise critical facts in real time.
By constructing a shared graph-vector database, newly spawned sub-agents
can rapidly inherit the team's cumulative knowledge, avoiding redundant
retrieval and cognitive divergence.

% ── 4.4  Reasoning Enhancement Mechanisms ─────────────────────────────────────
\subsection{Reasoning Enhancement Mechanisms}
\label{sec:reasoning}

\subsubsection{Collaborative Chain-of-Thought}
\label{sec:cot}

Chain-of-Thought (CoT) prompting \citep{wei2022chain_of_thought} has
greatly unlocked the reasoning potential of LLMs.
Self-Consistency \citep{wang2023self_consistency} samples multiple reasoning
paths from one model and selects an answer by marginalising over those paths.
Multi-agent systems adapt this diversity principle by assigning independent
reasoners and then using debate or confidence-aware aggregation
\citep{du2024multiagent_debate,chen2024reconcile}.

Each agent independently samples $K$ reasoning paths
$\{(r_k, a_k)\}_{k=1}^{K}$ for the same question, and majority voting
selects the final answer:
\begin{equation}
  \hat{a} = \underset{a\in\mathcal{Y}}{\arg\max}\sum_{k=1}^{K}
  \mathbf{1}\!\left[f_\theta(x,r_k)=a\right]
  \label{eq:self_consistency}
\end{equation}
where $f_\theta(x,r_k)$ is the answer produced by the $k$-th reasoning
chain and $\mathbf{1}[\cdot]$ is the indicator function.
The equation describes the single-model Self-Consistency baseline; exchange
between agents is an additional collaboration design, not a property of the
original method.

\subsubsection{Chain-of-Agents Reasoning}
\label{sec:chain_of_agents}

For ultra-long-text QA, the attention mechanism of conventional
Transformers faces severe computational bottlenecks.
Chain-of-Agents \citep{zhang2024chain_agents}, Tree of Thoughts
\citep{yao2023tree_of_thoughts}, and Graph of Thoughts
\citep{besta2024graph_of_thoughts} distribute reasoning across multiple
serially or graph-connected agents.
In Chain-of-Agents, worker agents process successive text chunks and pass
forward condensed evidence before a manager produces the answer. LONGAGENT
\citep{zhao2024longagent} uses a different leader--member design in which a
leader coordinates member agents over long-document segments. Tree- and
graph-structured reasoning provide related search topologies, but are not
necessarily multi-agent systems in their original formulations.

\subsubsection{Self-Reflection and Mutual-Critique Mechanisms}
\label{sec:reflexion}

To suppress hallucination, agents must be capable of iterative
self-examination.
Reflexion \citep{shinn2023reflexion} stores linguistic feedback from prior
trials in episodic memory, whereas Self-Refine
\citep{madaan2023self_refine} iterates between generation, feedback, and
revision without additional model training.
In addition, combined with Chain-of-Verification
\citep{dhuliawala2023chain_of_verification}, the system can configure
dedicated \emph{Critic} agents that rigorously audit generated results for
safety, factuality, and logical coherence through mutual-critique mechanisms.

% ── 4.5  Framework Comparison ─────────────────────────────────────────────────
\subsection{Open-source Frameworks and Toolkits}
\label{sec:framework_compare}

Since 2023, the open-source community has released numerous multi-agent
development frameworks.
Table~\ref{tab:framework_compare} provides a comprehensive comparison of
eight representative frameworks or orchestration toolkits across architecture type, communication
mechanism, task decomposition strategy, memory design, and typical QA
application domain.

\begin{table*}[pos=t]
  \centering
  \caption{Technical implementation comparison of representative multi-agent
           collaborative frameworks.}
  \label{tab:framework_compare}
  \resizebox{\linewidth}{!}{%
  \begin{tabular}{l c c c l l c l}
    \toprule
    \textbf{Framework} & \textbf{Year} & \textbf{Institution} &
    \textbf{Architecture} & \textbf{Communication} &
    \textbf{Task Decomposition} & \textbf{Memory} &
    \textbf{Typical QA Application} \\
    \midrule
    CAMEL \citep{li2023camel} & 2023 & KAUST et al. & Flat &
      Prompt injection, role-play dialogue &
      Role-based spontaneous decomposition &
      Short-term context dialogue tree &
      Exploratory dialogue, general task simulation \\
    ChatDev \citep{qian2023chatdev} & 2023 & Tsinghua Univ. & Hierarchical &
      Role-based dialogue along a Chat Chain &
      Waterfall-model static decomposition &
      Local code files and documents &
      Software development loop, code generation \\
    MetaGPT \citep{hong2024metagpt} & 2023 & DeepWisdom & Hierarchical &
      Structured SOP messages &
      Explicit role and responsibility division &
      Message queue and project documents &
      Enterprise software engineering, complex workflows \\
    AutoGen \citep{wu2024autogen} & 2023 & Microsoft & Configurable &
      Conversable-agent messages; multiple conversation patterns &
      Developer-defined orchestration and speaker selection &
      Conversation context + optional external tools &
      Complex conversational reasoning, code debugging \\
    AgentVerse \citep{chen2024agentverse} & 2024 & Tsinghua et al. & Hybrid &
      Explicit messages + debate protocol &
      Dynamic expert recruitment and phase switching &
      Global blackboard memory &
      NLP evaluation, multi-step logical problem-solving \\
    DSPy \citep{khattab2024dspy} & 2024 & Stanford & Modular pipeline toolkit &
      Modular declarative function calls &
      Metric-driven compilation of prompts and demonstrations &
      Module state and optimised demonstrations &
      QA and RAG pipeline optimisation \\
    LangGraph \citep{duan2024langgraph_crewai} & 2024 & LangChain & Hybrid &
      State-graph edge transitions &
      Developer-defined flow nodes &
      Persistent state-graph storage with checkpoint recovery &
      Production-grade fault-tolerant long-running QA \\
    GPTSwarm \citep{zhuge2024gptswarm} & 2024 & Swiss AI Lab & Hybrid &
      Computational graph topology connections &
      Automatic optimal graph-structure search &
      Intra-node state maintenance &
      High-performance optimisation, large-scale parallel inference \\
    \bottomrule
  \end{tabular}}
\end{table*}

Three clear evolutionary trends emerge from the cross-framework comparison.
\emph{First}, \textbf{architecture and scheduling are shifting from rigid
to dynamically flexible}: early frameworks such as ChatDev adopted stiff
pipeline structures, while later ones such as AgentVerse and GPTSwarm
support dynamic graph transitions and expert recruitment.
\emph{Second}, \textbf{communication protocols are increasingly
standardised}: developers are abandoning plain-text dialogue in favour of
highly structured JSON messages or state-graph transmission, which paves
the way for the adoption of protocols such as MCP
\citep{anthropic2024mcp}.
\emph{Third}, \textbf{memory and external tooling are becoming
first-class citizens}: leading frameworks now abstract RAG retrieval and
external tool invocation as built-in system components, substantially
lowering the development barrier for multi-agent QA systems.
